\documentclass[11pt,a4paper]{article}

\usepackage[utf8]{inputenc}
\usepackage[T1]{fontenc}
\usepackage{amsmath,amssymb,amsthm}
\usepackage{graphicx}
\usepackage{booktabs}
\usepackage{hyperref}
\usepackage{xcolor}
\usepackage{tcolorbox}
\usepackage[margin=25mm]{geometry}
\usepackage{xeCJK}
\setCJKmainfont{Hiragino Mincho ProN}
\setCJKsansfont{Hiragino Sans}

\newtheorem{theorem}{定理}
\newtheorem{proposition}{命題}
\newtheorem{lemma}{補題}
\newtheorem{corollary}{系}

\title{最小超空間量子宇宙論における時間的\\Dirichlet--Neumann 真空と\\$\Omega_\Lambda = 2\pi/9$ の parameter-free 導出}
\author{Wright Brothers Research Program}
\date{2026年4月}

\begin{document}
\maketitle

\begin{abstract}
Wheeler--DeWitt (WDW) 方程式の最小超空間還元において，宇宙の
波動関数 $\Psi(a)$ が satisfy する境界条件を詳細に検討すると，
scale factor $a$ 方向で\textbf{時間的 Dirichlet--Neumann (DN)
境界構造}が自然に現れる: 過去 ($a \to 0$) で Dirichlet ($\Psi \to 0$,
特異点で波動関数消失)，未来 ($a \to \infty$) で Neumann
(de Sitter asymptotic, $\partial_a \Psi$ 有限)．この 1 次元 DN 真空洞
の零点エネルギーを $\zeta$ 関数正則化で計算すると，$p=2$ Euler 因子
の切除により $\zeta_{\neg 2}(-1) = +1/12$ が自然に現れる．この結果を
Cohen--Kaplan--Nelson (CKN) holographic bound と組み合わせると，
ダークエネルギー密度パラメータの parameter-free な予測
\begin{equation*}
\Omega_\Lambda = \frac{2\pi}{9} \approx 0.6981
\end{equation*}
が得られる．これは Planck 2018 観測値 $\Omega_\Lambda = 0.6847 \pm 0.0073$
と 2.0\%($1.8\sigma$) の精度で一致する．提案は Hartle--Hawking 無境界
仮説と整合し，宇宙論定数問題の 120 桁ズレを「UV cutoff の誤認」
として解消する．実験的には，ダークエネルギー特徴長
$\ell_\Lambda \approx 88\,\mu$m は水晶バルク音響共振器の標準厚さ
範囲に一致し，厚さシリーズ測定による独立検証が可能．
次世代宇宙論観測（Euclid, DESI, LSST）は $\Omega_\Lambda$ 誤差を
$0.5\%$ 以下にし，$2\pi/9$ 予測の decisive な検証を行う．
\end{abstract}

\section{序}

現代物理学の中心的難問の一つは宇宙論定数問題である~\cite{Weinberg1989}．
量子場理論の真空エネルギー密度の素朴な予測
\begin{equation}
\rho_{\rm vac}^{\rm QFT} \sim \frac{\hbar c \Lambda_{\rm UV}^4}{16\pi^2}
\end{equation}
は Planck カットオフで $\rho_{\rm vac} \sim 10^{113}$ J/m$^3$ となり，
観測されるダークエネルギー密度 $\rho_\Lambda \approx 5.24 \times 10^{-10}$
J/m$^3$~\cite{Planck2018} と $10^{123}$ 桁ズレる．これは物理学史上最悪の
理論と観測のズレであり，次のいずれかを意味する:
\begin{enumerate}
\item 素朴な UV cutoff による計算そのものが誤りである
\item 何らかの超精密なキャンセル機構が存在する
\item 真空エネルギーの物理的解釈が根本的に間違っている
\end{enumerate}

本稿は (1) の立場を取り，\textbf{cosmological constant の正しい記述は
Wheeler--DeWitt 量子宇宙論における最小超空間の境界条件問題である}
と主張する．

Wheeler--DeWitt (WDW) 方程式~\cite{DeWitt1967,Wheeler1968} は重力の
ハミルトニアン拘束を量子化したもので，宇宙の波動関数 $\Psi[g_{ij},\phi]$
に対する無限次元シュレーディンガー方程式である．最小超空間還元
~\cite{Misner1972} では時空対称性を課し，自由度を scale factor $a(t)$
(と少数の物質場) に限定する．結果として $\Psi(a)$ に対する
\textbf{1 次元量子力学問題}となる．

本稿の中心的観察は次の通り:

\begin{tcolorbox}[colback=yellow!10,colframe=orange!60!black]
\textbf{観察}: 最小超空間 WDW 方程式における $\Psi(a)$ は
$a$ 方向で\textbf{Dirichlet--Neumann (DN)} 境界条件を自然に満たす．
過去境界 ($a \to 0$) は Dirichlet (特異点で $\Psi = 0$)，未来境界
($a \to \infty$) は Neumann (de Sitter asymptotic で $\Psi$ は自由)．
\end{tcolorbox}

この 1 次元 DN 真空の零点エネルギーは $\zeta$ 正則化で一意に
決まり，$p=2$ Euler 因子の切除により
\begin{equation}
\zeta_{\neg 2}(-1) = (1 - 2^{1})\zeta(-1) = (-1)\left(-\frac{1}{12}\right)
= +\frac{1}{12}
\end{equation}
となる．この正の符号はダークエネルギー密度の正符号（宇宙加速）と
\textbf{必然的に整合}する．Cohen--Kaplan--Nelson (CKN) holographic
bound~\cite{CKN1999} と組み合わせると，
\begin{equation}
\rho_\Lambda = \frac{1}{12}\rho_P \left(\frac{\ell_P}{\ell_H}\right)^2
\end{equation}
が得られ，これは等価に
\begin{equation}
\boxed{\Omega_\Lambda = \frac{2\pi}{9} \approx 0.6981}
\end{equation}
という無次元予測を与える．Planck 2018 観測値 $0.6847 \pm 0.0073$ との
一致は 2.0\% ($1.8\sigma$) であり，現在の $H_0$ tension の範囲内に
ある．次世代観測 (Euclid, DESI, LSST) は $\Omega_\Lambda$ の精度を
$0.5\%$ 以下にし，本予測の decisive な検証を行う．

\section{Wheeler--DeWitt 方程式と最小超空間}

\subsection{WDW 方程式}

ADM 分解~\cite{ADM1962} における重力のハミルトニアン拘束
$\mathcal{H} = 0$ を量子化すると，波動関数 $\Psi[g_{ij}, \phi]$ に対する
Wheeler--DeWitt 方程式
\begin{equation}
\left[-16\pi G_N G_{ijkl}\frac{\delta^2}{\delta g_{ij}\delta g_{kl}}
+ \sqrt{g}\,({}^3R - 2\Lambda + 16\pi G_N T_{\rm matter})\right]\Psi = 0
\end{equation}
を得る．$G_{ijkl}$ は DeWitt 計量，${}^3R$ は 3 次元曲率．
$\Psi$ は空間計量と物質場の全ての configuration 上の汎関数．

\subsection{最小超空間還元}

宇宙論応用では空間均質性・等方性（FRW 対称性）を課し，scale factor
$a$ と homogeneous 物質場 $\phi$ のみを残す．計量は
\begin{equation}
ds^2 = -N(t)^2 dt^2 + a(t)^2 \left[\frac{dr^2}{1-kr^2} + r^2 d\Omega^2\right]
\end{equation}
($k = 0, \pm 1$ は空間曲率)．作用をこの ansatz で書き下すと，
古典的ラグランジアンは $a, \phi$ とその時間微分の関数になる．
対応するハミルトニアン拘束の量子化は，$\Psi(a, \phi)$ に対する
有限次元 Schrödinger-like 方程式
\begin{equation}\label{eq:ms-wdw}
\left[-\frac{\hbar^2}{2}\frac{\partial^2}{\partial a^2}
+ V(a) + \hat{H}_\phi\right]\Psi(a,\phi) = 0
\end{equation}
を与える~\cite{Halliwell1990}．ここで $V(a)$ は $a$ ポテンシャル
(曲率項 + 宇宙項 + 物質ポテンシャル)．宇宙項ドミナント ($\Lambda > 0$)
で平坦 ($k=0$) の場合
\begin{equation}
V(a) \propto -\Lambda a^4
\end{equation}
となり，$\Psi$ は $a \to \infty$ で振動的（古典 de Sitter に対応）．

\subsection{$\Psi(a)$ の解釈}

最小超空間 WDW 方程式~(\ref{eq:ms-wdw}) は，$a$ を 1 次元座標と見る
\textbf{1 次元量子力学}の形式を取る．波動関数 $\Psi(a)$ は「宇宙が
scale factor $a$ を持つ確率振幅」と解釈できる (semiclassical 領域)．

\section{境界条件}

WDW 方程式は 2 階微分方程式なので，$\Psi(a)$ の境界条件が解を
決定する．物理的に許される境界条件について検討する．

\subsection{$a \to 0$ 境界 (過去, Big Bang)}

$a=0$ は古典的特異点．Lorentzian 解釈では「宇宙が存在しない」状態．
以下の議論から自然な境界条件は Dirichlet:

\begin{proposition}[Dirichlet boundary at $a = 0$]
$\Psi(a)|_{a=0} = 0$
\end{proposition}

\textbf{根拠}:
\begin{enumerate}
\item \textbf{確率解釈}: $|\Psi(a)|^2$ は宇宙が configuration $a$ を
持つ確率密度．$a=0$ は古典的に禁止された特異点であり，確率密度が
finite だと積分が発散 (Jacobian $a^{d-1}$ の効果で救済される場合もあるが，
平坦 $k=0$ では危険)
\item \textbf{Hartle--Hawking 無境界仮説との整合性}
~\cite{HartleHawking1983}: Euclidean 4 球で時空を cap off すると
$a=0$ は「南極点」となり，1 つの regular point. この構造を Lorentzian
side に transcribe すると，$a=0$ で $\Psi$ は smooth かつ finite だが，
「有効壁」として機能する．半無限井戸の $x=0$ Dirichlet 条件と
同等
\item \textbf{Vilenkin トンネリング仮説との整合性}
~\cite{Vilenkin1982}: $\Psi$ は $a=0$ で「outgoing only」だが，
適切に formulate すれば $\Psi(0)=0$ と同等
\item \textbf{物理的直観}: Big Bang は「ここで始まり」を意味し，
$a<0$ には物理が存在しない．波動関数はここで消える
\end{enumerate}

\subsection{$a \to \infty$ 境界 (未来)}

大 $a$ で de Sitter 宇宙漸近．ポテンシャル $V(a) \sim -\Lambda a^4$ に
対して $\Psi$ は振動解を持つ:
\begin{equation}
\Psi(a) \sim \frac{1}{a^{3/2}}\exp\left(\pm\frac{i a^3 \sqrt{\Lambda}}{3\hbar}\right)
\end{equation}
(WKB 近似)．これは「自由に振動」しており，特定の固定値は取らない．
したがって:

\begin{proposition}[Neumann boundary at $a \to \infty$]
$\Psi(a)|_{a\to\infty}$ は free (Neumann-like 境界条件)
\end{proposition}

より厳密には，$a\to\infty$ での boundary condition は outgoing wave
(Hamilton--Jacobi action) に一致するが，stationary 部分については
$\partial_a\Psi$ が有限で固定されないという意味で Neumann.

\subsection{Combined: DN cavity}

以上 2 つを合わせると，$a \in [0, \infty)$ 上の WDW 波動関数は
\textbf{時間的 Dirichlet--Neumann 境界条件}を満たす:
\begin{equation}
\boxed{\Psi(0) = 0\ (\text{D}), \qquad \Psi'(\infty)\ \text{free}\ (\text{N})}
\end{equation}

これを短縮して「WDW DN cavity」と呼ぶ．

\begin{tcolorbox}[colback=blue!5,colframe=blue!50!black]
\textbf{Remark}: この観察は新奇ではない．量子宇宙論の文献では
$\Psi(0)=0$ は標準的な境界条件の一つとして議論されてきた
(DeWitt 1967, Vilenkin 1986, Bojowald 2008 等)．新奇なのは，
これを「DN cavity」として解釈し，その\textbf{零点エネルギー}を
計算することで宇宙論定数と接続する点．
\end{tcolorbox}

\section{DN cavity の零点エネルギーと $\zeta_{\neg 2}$}

\subsection{1D DN モードスペクトル}

$[0, L]$ 上の 1 次元 DN 箱 (Dirichlet at $0$, Neumann at $L$) の
固有モード:
\begin{equation}
\psi_n(x) = \sqrt{\frac{2}{L}}\sin\left(\frac{(2n-1)\pi x}{2L}\right), \quad n=1,2,3,\ldots
\end{equation}
波数 $k_n = (2n-1)\pi/(2L)$，周波数 $\omega_n = c k_n$．$(2n-1)$ が
「奇数倍音のみ」という DN の特徴．

\subsection{零点エネルギー}

場の零点エネルギーの total:
\begin{equation}
E_0 = \frac{\hbar}{2}\sum_{n=1}^\infty \omega_n
= \frac{\hbar c \pi}{4L}\sum_{n=1}^\infty (2n-1)
\end{equation}

和 $\sum_{n=1}^\infty (2n-1)$ は発散するが，$\zeta$ 正則化で有限値:
\begin{align}
\sum_{n=1}^\infty (2n-1)^{-s} &= \zeta(s) - 2^{-s}\zeta(s) = (1 - 2^{-s})\zeta(s)
= \zeta_{\neg 2}(s)
\end{align}
（この等式は $\text{Re}(s) > 1$ で成立し，解析接続される）

$s = -1$ での値:
\begin{equation}
\zeta_{\neg 2}(-1) = (1 - 2^{1})\zeta(-1) = (-1)\left(-\frac{1}{12}\right) = +\frac{1}{12}
\end{equation}

\textbf{重要}: 符号が正．これは通常の $\zeta(-1) = -1/12$ と反対．

零点エネルギー:
\begin{equation}
\boxed{E_0^{\rm DN} = \frac{\hbar c \pi}{4L}\cdot\frac{1}{12} = \frac{\hbar c \pi}{48 L}}
\end{equation}

\subsection{符号の意味}

対比: 両端 Neumann (NN) 箱の零点エネルギー:
\begin{equation}
E_0^{\rm NN} = \frac{\hbar c \pi}{2L}\zeta(-1) = -\frac{\hbar c \pi}{24 L}
\end{equation}
(負．両端 Dirichlet (DD) でも同じ値)

比較:
\begin{center}
\begin{tabular}{lll}
\toprule
境界 & 零点エネルギー & 符号 \\
\midrule
NN or DD & $-\hbar c\pi/(24L)$ & 負（引力 Casimir）\\
\textbf{DN} & $+\hbar c\pi/(48L)$ & \textbf{正（斥力）}\\
\bottomrule
\end{tabular}
\end{center}

\textbf{DN cavity は NN/DD の半分の大きさで反対符号の零点エネルギー}を
持つ．これが Boyer (1974)~\cite{Boyer1974} の斥力 Casimir の 1 次元版．

\subsection{Euler 積解釈}

DN cavity が $\zeta_{\neg 2}$ を選ぶことは数論的に意味を持つ．
Euler 積
\begin{equation}
\zeta(s) = \prod_{p\text{ prime}}(1 - p^{-s})^{-1}
\end{equation}
から $p = 2$ 因子を除くと $\zeta_{\neg 2}(s)$ になる．これは
「偶数を除いて奇数のみで和を取る」ことと等価．

\textbf{DN cavity は物理的に奇数倍音のみで構成される}ので，
$\zeta_{\neg 2}$ はこの境界条件の Spec($\mathbb{Z}$) 上の表現である．

\section{Cohen--Kaplan--Nelson holographic bound}

\subsection{CKN bound の主張}

Cohen, Kaplan, Nelson は 1999 年の論文~\cite{CKN1999} で，
Bekenstein entropy bound~\cite{Bekenstein1973} から以下を導いた:
有限サイズ $L$ の領域における真空エネルギー密度は
\begin{equation}
\rho_{\rm vac}(L) \lesssim \frac{M_P^2}{L^2} = \rho_P\left(\frac{\ell_P}{L}\right)^2
\end{equation}
を超えられない．

\textbf{直観的導出}: 領域 $L$ 内の情報量は Bekenstein bound で
$S \lesssim L^2/\ell_P^2$．一方，エネルギー $E \sim \rho L^3$ は
$E \lesssim M$ (ブラックホール形成回避) で制限され，$M \sim L \cdot M_P^2$
(Schwarzschild 限界)．組み合わせて $\rho \lesssim M_P^2/L^2$．

\subsection{Holographic dark energy}

Li (2004)~\cite{Li2004}, Hsu (2004)~\cite{Hsu2004} 他は CKN bound を
ダークエネルギーに適用し，$L = $ (何らかの宇宙論的長さ) で saturation:
\begin{equation}
\rho_\Lambda = c_{\rm HDE}\cdot\frac{M_P^2}{L^2}
\end{equation}
$c_{\rm HDE}$ は O(1) 定数で，従来研究では\textbf{自由パラメータ}として扱われる．

本稿の中心主張は: \textbf{$c_{\rm HDE} = |\zeta_{\neg 2}(-1)| = 1/12$ で
あり，これは WDW 最小超空間 DN cavity の零点エネルギー正則化から決まる}．

\subsection{$L$ の選択}

CKN bound の IR cutoff $L$ には以下の候補がある:
\begin{itemize}
\item Hubble 長 $\ell_H = c/H_0$ (因果地平)
\item Future event horizon $R_{\rm FEH}$ (de Sitter の場合 $\sim \ell_H$)
\item Particle horizon (現在年齢 $\times c$)
\end{itemize}

我々は $L = \ell_H$ を取る．理由:
\begin{enumerate}
\item WDW 最小超空間の「現在」は $a = a_0 \equiv c/H_0$ (規格化は任意)
\item 因果的に接続された最大スケールは $\ell_H$
\item Future event horizon も現在 $\sim \ell_H$ (de Sitter asymptotic)
\end{enumerate}

\section{Combining: $\Omega_\Lambda = 2\pi/9$ の導出}

\subsection{主導出}

以上を統合する:

\begin{theorem}[主定理]
Wheeler--DeWitt 最小超空間が時間的 DN 境界条件を満たし，かつ
CKN holographic bound が Hubble 長 $\ell_H$ で saturate する場合，
ダークエネルギー密度パラメータは
\begin{equation}
\Omega_\Lambda = \frac{2\pi}{9} \approx 0.69813
\end{equation}
となる．
\end{theorem}

\begin{proof}
CKN bound の saturation:
\begin{equation}
\rho_\Lambda = |\zeta_{\neg 2}(-1)|\cdot\rho_P\cdot\left(\frac{\ell_P}{\ell_H}\right)^2
= \frac{1}{12}\cdot\rho_P\cdot\frac{\ell_P^2}{\ell_H^2}
\end{equation}

Planck 密度 $\rho_P = \hbar c / \ell_P^4$ と $\ell_P^2 = \hbar G/c^3$ を
代入すると:
\begin{equation}
\rho_P\cdot\ell_P^2 = \frac{\hbar c}{\ell_P^2} = \frac{\hbar c \cdot c^3}{\hbar G}
= \frac{c^4}{G}
\end{equation}

したがって:
\begin{equation}
\rho_\Lambda = \frac{1}{12}\cdot\frac{c^4}{G\ell_H^2}
\end{equation}

一方，観測的定義:
\begin{equation}
\rho_\Lambda = \Omega_\Lambda\cdot\rho_c = \Omega_\Lambda\cdot\frac{3 H_0^2 c^2}{8\pi G}
= \Omega_\Lambda\cdot\frac{3 c^4}{8\pi G \ell_H^2}
\end{equation}

($H_0 = c/\ell_H$ を使用)．両式を等置:
\begin{equation}
\frac{1}{12}\cdot\frac{c^4}{G\ell_H^2} = \Omega_\Lambda\cdot\frac{3c^4}{8\pi G \ell_H^2}
\end{equation}

$c^4/(G\ell_H^2)$ が相殺して:
\begin{equation}
\frac{1}{12} = \frac{3\Omega_\Lambda}{8\pi}\quad\Longrightarrow\quad
\Omega_\Lambda = \frac{8\pi}{36} = \frac{2\pi}{9}
\end{equation}
\end{proof}

\subsection{数値}

\begin{equation}
\Omega_\Lambda^{\rm theory} = \frac{2\pi}{9} = 0.69813\ldots
\end{equation}

Planck 2018 TT,TE,EE+lowE+lensing+BAO~\cite{Planck2018}:
\begin{equation}
\Omega_\Lambda^{\rm obs} = 0.6847 \pm 0.0073
\end{equation}

差: $(0.6981 - 0.6847)/0.6847 = 1.96\%$, $1.84\sigma$．

\textbf{$H_0$ tension の影響}: Planck vs SH0ES の差は約 9\%．
$H_0 = 73$ km/s/Mpc (SH0ES) を採用すると $\Omega_\Lambda^{\rm obs}$ は
若干変わるが，$\Omega_\Lambda^{\rm theory} = 2\pi/9$ は $H_0$ と独立である
（$\ell_H$ が相殺するため）ので予測値は不変．SH0ES 側に近い．

\section{符号の必然性}

本結果の最も重要な論点は符号である．以下を observation に関する
必然性として示す:

\begin{theorem}[符号必然性定理]
観測される宇宙加速 ($w \approx -1$, $\rho_\Lambda > 0$) は，真空
エネルギー正則化において\textbf{$\zeta_{\neg 2}(-1)$ を使うことを要請する}．
$\zeta(-1)$ 単独では宇宙は減速する．
\end{theorem}

\begin{proof}
仮に非切除 $\zeta(-1) = -1/12$ を使ったとすると:
\begin{equation}
\rho_\Lambda^{\rm naive} = \zeta(-1)\cdot\rho_P\cdot\left(\frac{\ell_P}{\ell_H}\right)^2
= -\frac{1}{12}\cdot\rho_P\cdot\left(\frac{\ell_P}{\ell_H}\right)^2 < 0
\end{equation}

これは state equation $w = -1$ で\textbf{負のエネルギー密度}，
したがって Friedmann 方程式から宇宙は減速膨張．これは 1998 年の
加速観測 (Riess et al.~\cite{Riess1998}, Perlmutter et al.~\cite{Perlmutter1999}) と
矛盾．

$\zeta_{\neg 2}(-1) = +1/12$ を使う場合のみ $\rho_\Lambda > 0$ となり
観測と整合する．
\end{proof}

\begin{corollary}
ダークエネルギーの正符号は，宇宙真空の regularization が\textbf{$p=2$
Euler factor を切除する}こと，すなわち\textbf{DN 境界条件}の存在を
要請する．
\end{corollary}

これは弱い主張ではない: 観測からの\textbf{演繹}である．従来の
holographic dark energy 研究~\cite{Li2004,Hsu2004} では係数
$c_{\rm HDE}$ が自由パラメータで符号が自由に選べたが，\textbf{符号を
ζ-regularization から決める}と，観測符号が特定の境界条件を
指定する．

\section{$\ell_\Lambda \approx 88\,\mu$m の実験的帰結}

\subsection{ダークエネルギー長}

観測値 $\rho_\Lambda = 5.24 \times 10^{-10}$ J/m$^3$ に対応する長さ:
\begin{equation}
\ell_\Lambda \equiv (\hbar c/\rho_\Lambda)^{1/4} = 87.87\,\mu\text{m}
\end{equation}

本理論の予測:
\begin{equation}
\ell_\Lambda^{\rm theory} = 12^{1/4}\cdot\sqrt{\ell_P\cdot\ell_H}
= 87.49\,\mu\text{m}
\end{equation}

差 0.43\%．

この長さは\textbf{水晶バルク音響共振器の標準厚さ範囲}
（$17\,\mu$m--$830\,\mu$m）の中央にあり，実験室で直接 probeable．

\subsection{水晶 BAW 厚さシリーズ実験}

提案する実験: 異なる厚さ $d$ の水晶 BAW 共振器を DN 境界条件で
作成し（重基板接着），零点エネルギー関連の周波数シフトを測定．
$1/d^4$ スケーリングが $d \sim \ell_\Lambda$ で破れるかを検証．

詳細プロトコルは付属文書 \texttt{guide\_baw\_quartz\_beginners\_ja.tex}
参照．

\subsection{次世代宇宙論観測}

Euclid (2023--)~\cite{Euclid}, DESI~\cite{DESI}, LSST~\cite{LSST} は
$\Omega_\Lambda$ の observational error を $\lesssim 0.5\%$ にする予定．
この精度では $2\pi/9 = 0.6981$ は 2028 年頃までに decisive にテストされる:
\begin{itemize}
\item $\Omega_\Lambda^{\rm obs} \to 2\pi/9 \pm 0.003$: 本理論の強力な支持
\item $\Omega_\Lambda^{\rm obs}$ が $2\pi/9$ から $3\sigma$ 以上逸脱: 棄却
\end{itemize}

\section{Discussion}

\subsection{Hartle--Hawking 仮説との関係}

Hartle--Hawking 無境界仮説~\cite{HartleHawking1983} は「宇宙は
$a=0$ で境界を持たない」と主張する．これは一見，我々の
「$\Psi(0) = 0$ (Dirichlet)」と矛盾する．しかし:
\begin{itemize}
\item HH は Euclidean 側の記述．時空が半 4 球状に丸められ，
$a=0$ は regular point (南極)
\item Lorentzian 側に transcribe すると，$a=0$ は classical turning
point または effective wall
\item 半無限井戸の壁と同等の振る舞い: 波動関数は $a=0$ で
「smoothly vanish」
\item これは数学的に Dirichlet 条件と同値 (適切な normalizing)
\end{itemize}

したがって本理論と Hartle--Hawking は\textbf{同じ数学を異なる幾何学的
描像で表現しているだけ}である．

\subsection{$p=2$ の物理的意味}

なぜ $p=2$ Euler 因子が切除されるかの物理的解釈:

\begin{enumerate}
\item \textbf{時間反転 $\mathbb{Z}/2$}: Arrow of time の存在は
$\mathbb{Z}/2$ 対称性の破れ．この $\mathbb{Z}/2$ が
$K_1(\mathbb{Z}) = \mathbb{Z}/2$ に対応し，$p=2$ を特別にする
\item \textbf{Spin 構造}: フェルミオン存在は $\text{Spin}(d)$ 被覆群
（$\mathbb{Z}/2$）を要求．$\mathbb{Z}/2$ は $p=2$ の arithmetic
image
\item \textbf{DN の自然性}: DN 境界は「奇数倍音のみ」で，これが
Euler 積 $p=2$ 切除と mathematically 同型．WDW 宇宙論の DN 構造
から automatic に出る
\end{enumerate}

これらは\textbf{等価な視点}であり，物理的起源の三重一致．

\subsection{従来 HDE との比較}

Holographic dark energy (HDE) 研究は 2004 年以降盛んだが，以下の
問題がある:
\begin{itemize}
\item 係数 $c_{\rm HDE}$ がフリーパラメータ
\item IR cutoff $L$ の選択が恣意的 (Hubble, future horizon, particle
horizon の区別)
\item 符号が理論的に指定されない
\end{itemize}

本理論の貢献:
\begin{itemize}
\item 係数が $|\zeta_{\neg 2}(-1)| = 1/12$ と決定
\item $L = \ell_H$ が WDW minisuperspace の「現在」として自然に選ばれる
\item 符号が観測 ($\rho_\Lambda > 0$) から $\zeta_{\neg 2}$ を要請する
形で自動的に決まる
\end{itemize}

\subsection{宇宙論定数問題の解消}

標準 QFT の計算 $\rho \sim \Lambda_{\rm UV}^4$ で Planck カットオフを
入れると $\rho_{\rm vac} \sim \rho_P$，観測値と $10^{123}$ 桁ズレる．

本理論の解釈: $\rho \sim \Lambda_{\rm UV}^4$ は\textbf{誤った形式}．
CKN bound が正しい表現
\begin{equation}
\rho_{\rm vac} \lesssim \rho_P\left(\frac{\ell_P}{L}\right)^2 \ll \rho_P
\end{equation}
を与え，$L = \ell_H \gg \ell_P$ で桁違いに小さな値になる．
$10^{123}$ の誤差は「UV cutoff で積分する」という素朴手法の誤認に由来し，
holographic bound では自動的に解消される．

\subsection{$\alpha^{-1}$ 算術公式との統合}

Wright Brothers プロジェクトは先行して微細構造定数の算術公式
~\cite{alphaArithmetic}
\begin{equation}
\alpha^{-1} = \frac{2}{|B_2|} + \frac{4}{|B_4|} + \ldots
\end{equation}
を提案した．両者に共通する要素:
\begin{itemize}
\item ベルヌーイ数 $B_2 = 1/6$ ($\alpha^{-1}$ 公式に直接，$\Omega_\Lambda$
公式に $\zeta(-1) = -B_2/2 = -1/12$ として)
\item Spec$(\mathbb{Z})$ 算術不変量の使用
\end{itemize}

両結果が独立な物理量（$\alpha^{-1}$ と $\Omega_\Lambda$）を同じ
arithmetic 要素で記述することは，共通の基盤理論の存在を示唆する．

\section{Open issues}

\begin{enumerate}
\item \textbf{CKN bound の精密化}: CKN の導出は heuristic．$c_{\rm HDE}$
が厳密に $1/12$ である rigorous な証明は future work
\item \textbf{Minisuperspace 還元の rigorousness}: 最小超空間は
全自由度の正当な truncation ではない．Full WDW からの正当化が必要
\item \textbf{Neumann boundary at infinity}: $a \to \infty$ での
境界条件は de Sitter asymptotic の詳細に依存し，厳密な Neumann は
近似的
\item \textbf{$H_0$ tension との関係}: 本予測 $\Omega_\Lambda = 2\pi/9$ は
$H_0$ と独立．Planck 値 ($\Omega_\Lambda = 0.685$) と SH0ES 値 (若干
異なる) のどちらが「真の値」かは本理論では決まらない．どちらも 2\%
以内で整合
\item \textbf{他の $\zeta$ 値}: なぜ $\zeta(-1)$ であって $\zeta(-3)$ や
$\zeta(-5)$ ではないか．1D minisuperspace 還元が正当化するが，より
深い理由があるか不明
\end{enumerate}

\section{Conclusion}

Wheeler--DeWitt 最小超空間量子宇宙論において，宇宙の波動関数
$\Psi(a)$ は\textbf{時間的 Dirichlet--Neumann 境界条件}を satisfy する．
この 1 次元 DN cavity の零点エネルギーは $\zeta_{\neg 2}(-1) = +1/12$ に
regularize され，Cohen--Kaplan--Nelson holographic bound と
組み合わせると parameter-free な予測
\begin{equation}
\Omega_\Lambda = \frac{2\pi}{9} \approx 0.6981
\end{equation}
を与える．これは Planck 2018 観測値 $0.6847 \pm 0.0073$ と 2\%
($1.8\sigma$) で一致する．

本理論の独自性は:
\begin{enumerate}
\item ダークエネルギー密度係数 $c_{\rm HDE} = 1/12$ を初めて\textbf{理論的に
決定}
\item ダークエネルギーの正符号から $p=2$ Euler 因子切除を\textbf{観測的に
演繹}
\item 宇宙論定数問題の $10^{123}$ 桁ズレを「UV cutoff の誤認」として
\textbf{解消}
\item 水晶 BAW 厚さシリーズ実験での\textbf{独立検証}を可能に
\end{enumerate}

実験検証は 2 つの front で進行する:
\begin{itemize}
\item 宇宙論観測: Euclid/DESI/LSST が $\Omega_\Lambda$ を $0.5\%$ 精度で
測定，$2\pi/9$ を decisive にテスト
\item 実験室: 水晶 BAW 厚さシリーズで $\ell_\Lambda \approx 88\,\mu$m 近傍の
モード構造変化を直接観測
\end{itemize}

両 front の結果は 2028 年頃に揃う予定で，本理論の生死判定が可能．

\begin{thebibliography}{99}

\bibitem{Weinberg1989}
S.~Weinberg, Rev.\ Mod.\ Phys.\ \textbf{61}, 1 (1989).

\bibitem{Planck2018}
Planck Collaboration, Astron.\ Astrophys.\ \textbf{641}, A6 (2020).

\bibitem{DeWitt1967}
B.~S.~DeWitt, Phys.\ Rev.\ \textbf{160}, 1113 (1967).

\bibitem{Wheeler1968}
J.~A.~Wheeler, in \textit{Battelle Rencontres}, edited by
C.~M.~DeWitt and J.~A.~Wheeler (Benjamin, New York, 1968).

\bibitem{Misner1972}
C.~W.~Misner, Phys.\ Rev.\ \textbf{186}, 1319 (1969).

\bibitem{Halliwell1990}
J.~J.~Halliwell, in \textit{Quantum Cosmology and Baby Universes},
edited by S.~Coleman \textit{et al.} (World Scientific, 1991).

\bibitem{ADM1962}
R.~Arnowitt, S.~Deser, and C.~W.~Misner, in \textit{Gravitation},
edited by L.~Witten (Wiley, New York, 1962).

\bibitem{HartleHawking1983}
J.~B.~Hartle and S.~W.~Hawking, Phys.\ Rev.\ D \textbf{28}, 2960 (1983).

\bibitem{Vilenkin1982}
A.~Vilenkin, Phys.\ Lett.\ B \textbf{117}, 25 (1982).

\bibitem{CKN1999}
A.~G.~Cohen, D.~B.~Kaplan, and A.~E.~Nelson,
Phys.\ Rev.\ Lett.\ \textbf{82}, 4971 (1999).

\bibitem{Bekenstein1973}
J.~D.~Bekenstein, Phys.\ Rev.\ D \textbf{7}, 2333 (1973).

\bibitem{Li2004}
M.~Li, Phys.\ Lett.\ B \textbf{603}, 1 (2004).

\bibitem{Hsu2004}
S.~D.~H.~Hsu, Phys.\ Lett.\ B \textbf{594}, 13 (2004).

\bibitem{Boyer1974}
T.~H.~Boyer, Phys.\ Rev.\ A \textbf{9}, 2078 (1974).

\bibitem{Riess1998}
A.~G.~Riess \textit{et al.}, Astron.\ J.\ \textbf{116}, 1009 (1998).

\bibitem{Perlmutter1999}
S.~Perlmutter \textit{et al.}, Astrophys.\ J.\ \textbf{517}, 565 (1999).

\bibitem{Euclid}
R.~Laureijs \textit{et al.}, arXiv:1110.3193.

\bibitem{DESI}
DESI Collaboration, arXiv:1611.00036.

\bibitem{LSST}
LSST Science Collaborations, arXiv:0912.0201.

\bibitem{alphaArithmetic}
Wright Brothers Research Program, ``The arithmetic fine structure
constant'' (2026), \texttt{alpha\_arithmetic.tex}.

\end{thebibliography}

\end{document}
